% test-009.tex -- tagged-PDF test del comando \spdist
%
% Compilar:  lualatex-dev test-009.tex
% Validar:   show-pdf-tags --xml test-009.pdf | rnv-wrapp
%            verapdf --flavour ua2 --format text test-009.pdf
%
% Cada caso es un parrafo numerado. Las lecturas "doc:" estan por escuchar
% con NVDA y MathCAT.
\DocumentMetadata
  {
    lang=es-CL,
    pdfversion=2.0,
    pdfstandard=ua-2,
    tagging=on,
    tagging-setup={math/setup={mathml-SE}},
  }
\documentclass{article}
\usepackage[spanish]{babel}
\usepackage{unicode-math}
\usepackage{spanish-mathml-intent}
\usepackage{hyperref}
\hypersetup
  {
    pdfdisplaydoctitle = true,
    pdftitle           = {test-009: el comando spdist},
  }

\begin{document}

\section*{A. El comando \texttt{\textbackslash spdist}}

% doc: (por escuchar)
1. Normal: $\spdist[var=X,dist=N]{\mu,\sigma^2}$.

% doc: (por escuchar)
2. Normal con desviación: $\spdist[var=X,dist=N]{\mu,\sigma}$.

% doc: (por escuchar)
3. Normal estándar: $\spdist[var=Z,dist=N]{0,1}$.

% doc: (por escuchar)
4. Normal con s: $\spdist[var=X,dist=N]{\bar{x},s}$.

% doc: (por escuchar)
5. Normal con s al cuadrado: $\spdist[var=X,dist=N]{\mu,s^2}$.

% doc: (por escuchar)
6. Normal con varianza numérica: $\spdist[var=X,dist=N,read-sty=variance]{5,4}$.

% doc: (por escuchar)
7. Normal sin variable: $\spdist[dist=N]{0,1}$.

% doc: (por escuchar)
8. Binomial: $\spdist[var=X,dist=B]{n,p}$.

% doc: (por escuchar)
9. Binomial con decimal: $\spdist[var=X,dist=B]{10;0,5}$.

% doc: (por escuchar)
10. Poisson: $\spdist[var=X,dist=P]{\lambda}$.

% doc: (por escuchar)
11. Bernoulli: $\spdist[var=X,dist=Be]{p}$.

% doc: (por escuchar)
12. Uniforme: $\spdist[var=X,dist=U]{a,b}$.

% doc: (por escuchar)
13. Uniforme con corchetes: $\spdist[var=X,dist=U,print-sty=brack]{a,b}$.

% doc: (por escuchar)
14. Uniforme discreta: $\spdist[var=X,dist=U,read-sty=discrete]{1,n}$.

% doc: (por escuchar)
15. Uniforme en un conjunto: $\spdist[var=X,dist=U,print-sty=set]{1,2,\dots,n}$.

% doc: (por escuchar)
16. Exponencial: $\spdist[var=X,dist=Exp]{\lambda}$.

% doc: (por escuchar)
17. Exponencial con tasa: $\spdist[var=X,dist=Exp,read-sty=rate]{\lambda}$.

% doc: (por escuchar)
18. Exponencial con media: $\spdist[var=X,dist=Exp,read-sty=mean]{\lambda}$.

% doc: (por escuchar)
19. Geométrica: $\spdist[var=X,dist=Ge]{p}$.

% doc: (por escuchar)
20. Hipergeométrica: $\spdist[var=X,dist=H]{N,n,p}$.

% doc: (por escuchar)
21. Gamma: $\spdist[var=X,dist=Gamma]{\alpha,\beta}$.

% doc: (por escuchar)
22. Gamma con tasa: $\spdist[var=X,dist=Gamma,read-sty=rate]{\alpha,\beta}$.

% doc: (por escuchar)
23. Gamma con escala: $\spdist[var=X,dist=Gamma,read-sty=scale]{k,\theta}$.

% doc: (por escuchar)
24. Beta: $\spdist[var=X,dist=Beta]{\alpha,\beta}$.

% doc: (por escuchar)
25. t de Student: $\spdist[var=X,dist=t]{n}$.

% doc: (por escuchar)
26. t de Student entre paréntesis: $\spdist[var=X,dist=t,print-sty=paren]{n}$.

% doc: (por escuchar)
27. Chi cuadrado: $\spdist[var=X,dist=chi]{n}$.

% doc: (por escuchar)
28. Chi cuadrado sin variable: $\spdist[dist=chi]{n}$.

% doc: (por escuchar)
29. F de Snedecor: $\spdist[var=X,dist=F]{m,n}$.

% doc: (por escuchar)
30. F de Snedecor entre paréntesis: $\spdist[var=X,dist=F,print-sty=paren]{m,n}$.

% doc: (por escuchar)
31. Alias Bi: $\spdist[var=X,dist=Bi]{n,p}$.

% doc: (por escuchar)
32. Alias Bin: $\spdist[var=X,dist=Bin]{n,p}$.

% doc: (por escuchar)
33. Alias Ber: $\spdist[var=X,dist=Ber]{p}$.

% doc: (por escuchar)
34. Alias Po: $\spdist[var=X,dist=Po]{\lambda}$.

% doc: (por escuchar)
35. Alias Geo: $\spdist[var=X,dist=Geo]{p}$.

% doc: (por escuchar)
36. Variable con subíndice: $\spdist[var=X_{1},dist=B]{n,p}$.

% doc: (por escuchar)
37. Varias variables: $\spdist[var={X_1,\dots,X_n},dist=N]{\mu,\sigma^2}$.

% doc: (por escuchar)
38. Lectura mathcat: $\spdist[var=X_1,dist=N,read-cmd=mathcat]{\mu,\sigma}$.

% doc: (por escuchar)
39. Silenciado: $\spdist[var=X,dist=N,silent=true]{\mu,\sigma^2}$.

% doc: (por escuchar)
40. Con read-arg: $\spdist[var=X,dist=N,read-arg={X sigue una distribución normal}]{\mu,\sigma^2}$.

% doc: (por escuchar)
41. Silenciado y read-arg: $\spdist[var=X,dist=N,silent=true,read-arg={nada}]{\mu,\sigma^2}$.

% doc: (por escuchar)
42. read-sty que no corresponde: $\spdist[var=Z,dist=N,read-sty=rate]{0,1}$.

% doc: (por escuchar)
43. Normal con desviación y exponente: $\spdist[var=X,dist=N,read-sty=std]{\mu,s^2}$.

\end{document}
